% =====================================================================
% DEPRECATED — kept for reference only.
% This file was the author's hand-written first-draft of the
% introduction. Its motivational opening (Dutch flood / 1953
% Watersnoodramp / Delta Works / wind storm risk) has been
% incorporated, with factual corrections, into Section
% sec:intro-motivation of thesis/introduction.tex during the
% P-Agent 2 production pass. The text below is preserved verbatim
% so that the author can recover original phrasing if desired, but
% it is NOT included in thesis-main.tex and should not be cited
% or referenced from any other thesis file.
%
% Reasons for not using this text directly:
%   1. Several factual claims need a citation or are imprecise
%      (e.g. "1000 kilometre wide stormfront", "10 out of 20 most
%      costly hazards", "return period of 1500 years for river
%      levels"). The Delta Commission's river norm is 1 250 years
%      (Davison, Padoan, Ribatet 2012, Section 2.2).
%   2. The phrase "wind speed" appears in the last sentence as the
%      main variable; the thesis is about wind gusts (FX) only.
%   3. The original mixes Dutch and English ("Sterker nog",
%      "achterland van de Hollandse IJssel", "Deltawerken",
%      "Afsluitdijk", "IJsselmeer") and contains the HTML entity
%      "&#x20;".
%   4. The current network size ("67 locations") is given as a
%      hard number; this is a data-lookup question marked as TODO
%      in thesis/introduction.tex.
% =====================================================================

% Original draft preserved below for reference only:
%
% The study of extreme weather events has become an increasingly
% important topic in econometrics. The amount of weather observations
% in both space and time resolution available has increased over the
% past decades by ...%. The Netherlands is situated in a river delta.
% The province south Holland is 3 meters below nap, with the lowest
% point 6.76 meters below nap, the achterland van de Hollandse
% IJssel. In 1953, the biggest natural disaster of the 20th century
% occurred, the Watersnoodramp. Dikes were breaking during springtij
% and a 1000 kilometer wide stormfront that approaches the Dutch
% coastal areas. The effect was a flood in parts of Zeeland and
% south holland, with 1836 dead. This led to the construction of the
% "Deltawerken". The commission responsible for construction used
% return period of extreme sea levels once every 10 000 years, and
% return periods once every 1500 years for river levels. Anywhere in
% the Netherlands, a radius of max 1 kilometer surrounds you with
% water. Logically, in the Netherlands the narrative on "natural
% disasters" is characterized by the fear of water. Wind plays an
% important role in these flood risk. Given the wind direction land
% inwards, and high wind speeds, the lake "IJsselmeer" can push
% large amounts of water into the IJssel, leading to river floods.
% The engineer Cornelis Lely made the Afsluitdijk to prevent the
% IJsselmeer of floods from the north sea back then. Therefore not
% only precipitation plays a role in the intensity of flood risk
% and storms, but also wind extremes. Sterker nog, 10 out of 20
% most costly (for the Dutch society) climate related hazards in
% the Netherlands, were windstorms. The Dutch society needs to not
% only be aware of flood risk, but also on windstorm risk. We will
% look whether there is time variation in the spatial dependence of
% wind speed at these KNMI locations. Since 1951, the KNMI has
% started recording weather observations at locations in the
% Netherlands, with a contemporary network of 67 locations spanning
% the up the North Sea down to Maastricht.
